\section*{\textbf{Research on Handwritten Digit Recognition Based on Convolutional Neural Networks}}


\section*{\textbf{Abstract}}
\addcontentsline{toc}{section}{Abstract}


Handwritten digit recognition is a classic problem in pattern recognition and image classification with wide applications. Traditional methods rely on hand-crafted features and exhibit limited generalization, whereas convolutional neural networks can automatically learn hierarchical features and are better suited to this task. Taking the MNIST dataset as the research object, this thesis designs a lightweight convolutional neural network: four $3\times3$ convolutional layers with batch normalization and ReLU, followed by two max-pooling operations and fully connected layers for classification, with about 0.47M parameters in total.


On this basis, the thesis investigates the roles of data augmentation, batch normalization, Dropout, and the Adam optimizer. Experiments show that the proposed method achieves a recognition accuracy of 99.32\% on the MNIST test set, outperforming Softmax regression, the multilayer perceptron, and LeNet-5, while the ablation study verifies the effectiveness of each training technique.


\vspace{\baselineskip}

{\noindent
\begingroup
\raggedright
\setbox0=\hbox{\textbf{Keywords}\quad}
\hangindent=\wd0
\hangafter=1
\makebox[\wd0][l]{\textbf{Keywords}\quad}convolutional~neural~network, handwritten~digit~recognition, MNIST, deep~learning, image~classification\par
\endgroup
}
